\subsection{Artículo especializado para publicar en revista indizada}
Bandeja y expediente del trámite.
\subsubsection*{Protocolo y aprobación del tema}
Valida o devuelve con observaciones cada documento; coteja el original cuando se indique.
\begin{longtable}{p{0.55\linewidth}p{0.35\linewidth}}
\toprule
\textbf{Documento} & \textbf{Observaciones}\\
\midrule
\endhead
Solicitud y protocolo firmados (pasante y asesor) & La plataforma lo genera. La plataforma lo genera con lo que escribiste; tu asesor lo firma por correo. Después lo revisa el comité de tu licenciatura.\\
Carta de presentación del pasante y asesor & Descárgala en blanco desde «Formatos en blanco». Solo para instituciones incorporadas.\\
\bottomrule
\end{longtable}

Aprobación del protocolo y oficio de aprobación del tema. Cuando el comité aprueba el protocolo, la plataforma genera el oficio de aprobación del tema; lo firman la Jefatura de Evaluación Profesional (Vo. Bo.) y la Subdirección Académica.
Datos que registras:
\begin{itemize}
  \item Fecha de aprobación del protocolo
  \item Observaciones del oficio (vacío: la leyenda del coasesor, si lo hay, o «Ninguna.»)
\end{itemize}
Se genera: Oficio de aprobación de tema y modalidad, que firma la Subdirección.
\subsubsection*{Desarrollo del trabajo escrito}
\subsubsection*{Solicitud de revisión del trabajo escrito (examen previo)}
Valida o devuelve con observaciones cada documento; coteja el original cuando se indique.
\begin{longtable}{p{0.55\linewidth}p{0.35\linewidth}}
\toprule
\textbf{Documento} & \textbf{Observaciones}\\
\midrule
\endhead
Solicitud de revisión firmada (pasante y asesor) & La plataforma lo genera. La firmas tú y tu asesor da su Vo. Bo. desde su correo.\\
Reporte de porcentaje de similitud (iThenticate) & No mayor al 20 \%.\\
Trabajo escrito en PDF para revisión & Con copia del oficio de aprobación de tema al inicio.\\
Solicitud de prórroga del tema (pasante y asesor) & La plataforma lo genera. Tu protocolo se aprobó hace más del plazo de exclusividad (Art. 86, REP). Captura avance y motivos en «Solicitudes y formatos adicionales», firma aquí y tu asesor da su Vo. Bo.; al validarla, la plataforma genera el oficio de aprobación de prórroga que firma la Subdirección.\\
\bottomrule
\end{longtable}

Fecha del examen previo y sínodo. Eliges fecha y hora y propones tu sínodo; el Coordinador de Docencia de tu licenciatura lo ratifica o modifica y firma desde el enlace de su correo. Después se generan los oficios individuales.
Se genera: Asignación de sínodo (Coordinación de Docencia), que firma la Subdirección.
\subsubsection*{Solicitud de sustentación y voto aprobatorio}
Valida o devuelve con observaciones cada documento; coteja el original cuando se indique.
\begin{longtable}{p{0.55\linewidth}p{0.35\linewidth}}
\toprule
\textbf{Documento} & \textbf{Observaciones}\\
\midrule
\endhead
Voto aprobatorio (lo firman tu asesor y tus revisores) & La plataforma lo genera. La plataforma lo genera y pide las firmas por correo.\\
Solicitud de sustentación firmada (pasante y asesor) & La plataforma lo genera. Se habilita cuando el voto aprobatorio está firmado.\\
\bottomrule
\end{longtable}

Autorización de impresión. 
Datos que registras:
\begin{itemize}
  \item Fecha de autorización
\end{itemize}
Se genera: Oficio de autorización de impresión, que firma la Subdirección.
\subsubsection*{Documentación para la evaluación profesional}
Valida o devuelve con observaciones cada documento; coteja el original cuando se indique.
\begin{longtable}{p{0.55\linewidth}p{0.35\linewidth}}
\toprule
\textbf{Documento} & \textbf{Observaciones}\\
\midrule
\endhead
Solicitud para presentar la evaluación profesional (trabajo escrito) & La plataforma lo genera. La plataforma lo genera con tus datos; lo firmas electrónicamente y la Subdirección lo autoriza.\\
Trabajo escrito final en PDF & Con oficio de autorización de impresión, resumen y carta de autorización para el Repositorio Institucional.\\
Reporte de porcentaje de similitud (iThenticate) & No mayor al 20 \%.\\
Constancia de pertenencia de la revista a los índices validados & Emitida por la Secretaría de Investigación y Estudios Avanzados (Art. 16, REP).\\
Carta de recepción o publicación del artículo expedida por la revista & Incluye criterios de publicación y documento de envío (Art. 20, REP).\\
Certificado de estudios de licenciatura & Descargado del Portal de alumnos.\\
Certificado de servicio social & \\
Certificado de estudios de bachillerato & \\
Constancia de autenticación de bachillerato & Se tramita en la escuela donde cursaste el bachillerato. Solo si tus estudios de bachillerato no fueron en instituciones dependientes o incorporadas a la UAEM.\\
Acta de nacimiento & Copia certificada vigente (del año en curso).\\
CURP & Debe mostrar la leyenda «Validada por el Registro Civil».\\
Constancia o certificado de aprobación del idioma inglés & Solo planes 98 en adelante, si en tu certificado no aparecen Inglés C1 e Inglés C2.\\
Constancia de no adeudo a la Biblioteca de la Facultad & El recibo de pago se tramita en Control Escolar de la Facultad; la constancia, en la Biblioteca.\\
Constancia de no adeudo a la Biblioteca Central & \\
Contra-recibo de la donación de libros a la Biblioteca de la Facultad & \\
Constancia de no adeudo a la Dirección de Control Escolar de la UAEM & Se solicita en el Portal de alumnos: controlescolar.uaemex.mx.\\
Recibo de no adeudo a la UAEM & Se solicita en el Portal de alumnos de Control Escolar.\\
Formato de datos personales (para el título) & La plataforma lo genera. La plataforma lo llena con tus datos; lo firmas electrónicamente.\\
Comprobante del sistema de Seguimiento de Egresados & Se obtiene en alumniversitario.uaemex.mx/1dae.php.\\
Dictamen de autorización de los HH. Consejos & La plataforma lo genera. Se solicita en la Subdirección Académica con el formato del Art. 5. Solo si excediste el plazo de 10 años para presentar la evaluación profesional (Art. 5, REP).\\
Carta de autorización o de excepción para publicar en el Repositorio Institucional & Entrega la que corresponda; descárgala en «Formatos en blanco».\\
Recibo de pago: derechos de evaluación profesional (\$3,024) & Se solicita en Control Escolar de la Facultad.\\
Recibo de pago: expedición de título (\$1,187) & \\
\bottomrule
\end{longtable}
\subsubsection*{Solicitud del examen definitivo}

Fecha del examen definitivo y jurado. Eliges fecha y hora y propones tu jurado; el Coordinador de Docencia lo ratifica o modifica y firma desde el enlace de su correo. Después se generan los oficios individuales.
Se genera: Asignación de jurado (Coordinación de Docencia), que firma la Subdirección.
\subsubsection*{Resultado de la evaluación}

Registrar resultado. 
Datos que registras:
\begin{itemize}
  \item Resultado
  \item Fecha de la evaluación
  \item Número de acta (de evaluación profesional; si se aplaza, de aplazamiento)
\end{itemize}
\subsubsection*{Verificación final del expediente}
Integración De Expediente: Artículo Especializado Publicado En Revista Arbitrada.

Documentos que verifica el responsable, en el orden del formato oficial:
\begin{enumerate}
  \item Dos fotografías tamaño semitítulo, blanco y negro, papel mate, fondo gris claro, de frente, con retoque, varones con camisa blanca, saco y corbata oscuros, mujeres con blazer oscuro --- \emph{físico}
  \item Acta de evaluación profesional --- \emph{físico}
  \item Copia del certificado de bachillerato con firma y sello de cotejo con el original --- \emph{plataforma}
  \item Constancia de autenticación si el bachillerato no es UAEM --- \emph{plataforma (sólo si aplica)}
  \item Dictamen de revalidación expedido por la Secretaría de Educación Pública (sólo si aplica) --- \emph{físico (sólo si aplica)}
  \item Acta de nacimiento --- \emph{plataforma}
  \item Certificado total de licenciatura --- \emph{plataforma}
  \item Solicitud por escrito del egresado dirigida a la Subdirección Académica con firma y sello de la misma --- \emph{plataforma}
  \item Carta de aceptación del documento en la revista arbitrada --- \emph{plataforma}
  \item Oficio de validación emitido por la Secretaría de Ciencia --- \emph{plataforma}
  \item Autorización para continuar con trámites de titulación con fundamento en el artículo 5° del Reglamento de Evaluación Profesional (sólo si aplica) --- \emph{plataforma (sólo si aplica)}
  \item Constancia de no adeudo a Servicios Escolares --- \emph{plataforma}
  \item Formato de revisión para autorización de evaluación profesional --- \emph{plataforma}
  \item Formato de integración de expediente --- \emph{plataforma}
  \item Formato de datos del interesado para expedición de título profesional --- \emph{plataforma}
\end{enumerate}
Marca con «Entregado y revisado» lo que se entrega en físico y con «No aplica» lo condicional; con todo cubierto, pulsa «Hago constar que el expediente está completo y verificado». Después, en \emph{Expedientes a Servicios Escolares}, elige los verificados y pulsa «Exportar carpeta comprimida para Servicios Escolares».
\subsubsection*{Pantallas del proceso}
\captura{articulo_titulacion_01_expediente_ronda00__1}{Expediente del trámite en la bandeja de Titulación. (parte 1 de 2)}
\captura{articulo_titulacion_01_expediente_ronda00__2}{Expediente del trámite en la bandeja de Titulación. (parte 2 de 2)}
\captura{articulo_titulacion_03_expedientes_bandeja}{Expediente del trámite en la bandeja de Titulación.}
\captura{articulo_titulacion_04_enviado}{Expedientes ya enviados a Servicios Escolares.}
\captura{articulo_titulacion_verificacion_01__1}{Lista de verificación final del expediente. (parte 1 de 2)}
\captura{articulo_titulacion_verificacion_01__2}{Lista de verificación final del expediente. (parte 2 de 2)}
\captura{articulo_titulacion_verificacion_02__1}{Expediente verificado y exportación para Servicios Escolares. (parte 1 de 2)}
\captura{articulo_titulacion_verificacion_02__2}{Expediente verificado y exportación para Servicios Escolares. (parte 2 de 2)}
